\exercisegroup

\begin{enumerate}
\def\labelenumi{\arabic{enumi}.}
\item
  Show that, in the set of orderings on a set E, the minimal elements (with respect to the order relation ``Γ is coarser than Γ'' between Γ and Γ'') are the total orderings on E, and that if Γ is any ordering on E, the graph of Γ is the intersection of the graphs of the total orderings on E which are coarser than Γ (apply Theorem 2 of no. 4). Deduce that every ordered set is isomorphic to a subset of a product of totally ordered sets.
\item
  Let E be an ordered set and let B be the set of subsets of E which are well-ordered by the induced ordering. Show that the relation ``X is a segment of Y'' on B is an order relation between X and Y and that B is inductive with respect to this order relation. Deduce that there exist well-ordered subsets of E which have no strict upper bound in E.
\item
  Let E be an ordered set. Show that there exist two subsets A, B of E such that \(A \cup B = E\) and \(A \cap B = \emptyset\) , and such that A is well-ordered and B has no least element (for example, take B to be the union of those subsets of E which have no least element). * Give an example in which there are several partitions of E into two subsets having these properties. *
\item
  An ordered set F is said to be partially well-ordered if every totally ordered subset of F is well-ordered. Show that in every ordered set E there exists a partially well-ordered subset which is cofinal in E. (Consider the set F of partially well-ordered subsets of E, and the order relation ``X ⊂ Y and no element of Y --- X is bounded above by any element of X'' between elements X and Y of F. Show that F is inductive with respect to this order relation.)
\item
  Let E be an ordered set and let J be the set of free subsets of E, ordered by the relation defined in §1, Exercise 5. Show that, if E is inductive, then J has a greatest element.
\end{enumerate}

¶ 6. Let E be an ordered set and let f be a mapping of E into E such that \(f(x) \geqslant x\) for all \(x \in E\) .

\begin{enumerate}
\def\labelenumi{(\alph{enumi})}
\item
  Let \(\mathfrak{G}\) be the set of subsets \(\mathbf{M}\) of \(\mathbf{E}\) with the following properties: (1) the relation \(x \in \mathbf{M}\) implies \(f(x) \in \mathbf{M}\) ; (2) if a non-empty subset of \(\mathbf{M}\) has a least upper bound in \(\mathbf{E}\) , then this least upper bound belongs to \(\mathbf{M}\) . For each \(a \in \mathbf{E}\) , show that the intersection \(\mathbf{C}_a\) of the sets of \(\mathfrak{G}\) which contain \(a\) also belongs to \(\mathfrak{G}\) ; that \(\mathbf{C}_a\) is well-ordered; and that if \(\mathbf{C}_a\) has an upper bound \(b\) in \(\mathbf{E}\) , then \(b \in \mathbf{C}_a\) and \(f(b) = b\) . \(\mathbf{C}_a\) is said to be the chain of \(a\) (with respect to the function \(f\) ). (Consider the set \(\mathfrak{M}\) whose elements are the empty set and the subsets \(\mathbf{X}\) of \(\mathbf{E}\) which contain \(a\) and have a least upper bound \(m\) in \(\mathbf{E}\) such that either \(m \notin \mathbf{X}\) or \(f(m) > m\) , any apply Lemma 3 of no. 3 to the set \(\mathfrak{M}\) .)
\item
  Deduce from (a) that if \(\mathbf{E}\) is inductive, then there exists \(b \in \mathbf{E}\) such that \(f(b) = b\) .
\end{enumerate}

\begin{enumerate}
\def\labelenumi{\arabic{enumi}.}
\setcounter{enumi}{6}
\tightlist
\item
  Let E be an ordered set and let F be the set of all closures (§ 1, Exercise 13) in E. Order F by putting \(u \leqslant v\) whenever \(u(x) \leqslant v(x)\) for all \(x \in E\) . Then F has a least element \(e\) , the identity mapping of E onto itself. For each \(u \in F\) let I(u) denote the set of elements of E which are invariant under \(u\) .
\end{enumerate}

\begin{enumerate}
\def\labelenumi{(\alph{enumi})}
\item
  Show that \(u \leqslant v\) in \(\mathbf{F}\) if and only if \(\mathbf{I}(v) \subset \mathbf{I}(u)\) .
\item
  Show that if every pair of elements of E has a greatest lower bound in E, then every pair of elements of F has a greatest lower bound in F. If E is a complete lattice, then so is F (§ 1, Exercise 11).
\item
  Show that if E is inductive (with respect to the relation \(\leqslant\) ), then every pair u, v of elements of F have a least upper bound in F. (Show that if \(f(x) = v(u(x))\) and if \(w(x)\) denotes the greatest element of the chain of x, relative to f (Exercise 6), then w is a closure in E and is the least upper bound of u and v.)
\end{enumerate}

¶ 8. An ordered set E is said to be ramified (on the right) if, for each pair of elements x, y of E such that x \textless{} y, there exists z \textgreater{} x such that y and z are not comparable. E is said to be completely ramified (on the right) if it is ramified and has no maximal elements. Every antidirected set (§ 1, Exercise 22) is ramified.

\begin{enumerate}
\def\labelenumi{(\alph{enumi})}
\item
  Let E be an ordered set and let a be an element of E. Let \(R_{a}\) denote the set of ramified subsets of E which have a as least element. Show that \(R_{a}\) , ordered by inclusion, has a maximal element.
\item
  If \(\mathbf{E}\) is branched (§ 1, Exercise 24), show that every maximal element of \(\mathfrak{R}_a\) is completely ramified.
\item
  Give an example of a branched set which is not ramified. The branched set defined in §1, Exercise 24 (c) is completely ramified.
\item
  Let E be a set in which each interval \(\leftarrow, x]\) is totally ordered. Show that E has an antidirected cofinal subset (§1, Exercise 22) (use (b)).
\end{enumerate}

\begin{enumerate}
\def\labelenumi{\arabic{enumi}.}
\setcounter{enumi}{8}
\item
  An ordinal sum \(\sum_{i\in I}E_i\) ( \(\S 1\) , Exercise 3) is well-ordered if and only if \(I\) and each \(E_i\) is well-ordered.
\item
  Let I be an ordered set and let \((\mathbf{E}_t)_{t \in I}\) be a family of ordered sets, all equal to the same ordered set E. Show that the ordinal sum \(\sum_{t \in I} \mathbf{E}_t\) if \(t \in I\) , then the order set is given by the sum of the orders \(\mathbf{E}_t\) (§ 1, Exercise 3) is isomorphic to the lexicographical product of the sequence \((\mathbf{F}_{\lambda})_{\lambda \in |\alpha, \beta|}\) , where the set \(\{\alpha, \beta\}\) of two distinct elements is well-ordered by the relation whose graph is \(\{(\alpha, \alpha), (\alpha, \beta), (\beta, \beta)\}\) , and where \(\mathbf{F}_{\alpha} = \mathbf{I}\) and \(\mathbf{F}_{\beta} = \mathbf{E}\) . This product is called the lexicographic product of \(\mathbf{E}\) by \(\mathbf{I}\) and is written E.I.
\end{enumerate}

¶ * 11. Let I be a well-ordered set and let \((E_t)_{t \in I}\) be a family of ordered sets, each of which contains at least two distinct comparable elements. Then the lexicographic product of the \(E_t\) is well-ordered if and only if each of the \(E_t\) is well-ordered and I is finite (if I is infinite, construct a strictly decreasing infinite sequence in the lexicographic product of the \(E_t\) ). *

\begin{enumerate}
\def\labelenumi{\arabic{enumi}.}
\setcounter{enumi}{11}
\tightlist
\item
  Let I be a totally ordered set and let \((\mathbf{E}_{t})_{t\in\mathbf{I}}\) be a family of ordered sets indexed by I. Let \(R\{x,y\}\) denote the following relation on \(E=\prod_{i\in I}E_{i}\) : ``the set of indices \(i\in I\) such that \(pr_{i}x\neq pr_{i}y\) is well-
\end{enumerate}

ordered, and if \(x\) is the least element of this subset of I, we have \(\mathrm{pr}_x x < \mathrm{pr}_x y\) . Show that \(R\{x, y\}\) is an order relation between \(x\) and \(y\) on E. If the \(E_t\) are totally ordered, show that the connected components of E with respect to the relation `` \(x\) and \(y\) are comparable'' (Chapter II, §6, Exercise 10) are totally ordered sets. Suppose that each \(E_t\) has at least two elements. Then E is totally ordered if and only if I is well-ordered and each \(E_t\) is totally ordered (use Exercise 3); and E is then the lexicographic product of the \(E_t\) .

\begin{enumerate}
\def\labelenumi{\arabic{enumi}.}
\setcounter{enumi}{12}
\tightlist
\item
  \begin{enumerate}
  \def\labelenumii{(\alph{enumii})}
  \tightlist
  \item
    Let \(\operatorname{Is}(\Gamma, \Gamma')\) be the relation `` \(\Gamma\) is an ordering (on E), and \(\Gamma'\) is an ordering (on E'), and there exists an isomorphism of E, ordered by \(\Gamma\) , onto E', ordered by \(\Gamma'\)''. Show that \(\operatorname{Is}(\Gamma, \Gamma')\) is an equivalence relation on every set whose elements are orderings. The term \(\tau_{\Delta}(\operatorname{Is}(\Gamma, \Delta))\) is an ordering called the order-type of \(\Gamma\) and is denoted by \(\operatorname{Ord}(\Gamma)\) , or \(\operatorname{Ord}(\mathbf{E})\) by abuse of notation. Two ordered sets are isomorphic if and only if their order-types are equal.
  \end{enumerate}
\end{enumerate}

\begin{enumerate}
\def\labelenumi{(\alph{enumi})}
\setcounter{enumi}{1}
\item
  Let \(\mathbb{R}\{\lambda, \mu\}\) be the relation: `` \(\lambda\) is an order-type, and \(\mu\) is an order-type, and there exists an isomorphism of the set ordered by \(\lambda\) onto a subset of the set ordered by \(\mu\) ''. Show that \(\mathbb{R}\{\lambda, \mu\}\) is a preorder relation between \(\lambda\) and \(\mu\) . It will be denoted by \(\lambda \prec \mu\) .
\item
  Let I be an ordered set and let \((\lambda_{t})_{t\in I}\) be a family of order-types indexed by I. The order-type of the ordinal sum (§ 1, Exercise 3) of the family of sets ordered by the \(\lambda_{t}(t\in I)\) is called the ordinal sum of the order-types \(\lambda_{t}(t\in I)\) and is denoted by \(\sum_{i\in I}\lambda_{i}\) . If \((E_{i})_{i\in I}\) is a family of ordered sets, the order-type of \(\sum_{i\in I}E_{i}\) is \(\sum_{i\in I}\operatorname{Ord}(E_{i})\) . If I is the ordinal sum of a family \((J_{x})_{x\in K}\) , show that
\end{enumerate}

\[
\sum_ {x \in K} \left(\sum_ {i \in J _ {x}} \lambda_ {i}\right) = \sum_ {i \in I} \lambda_ {i}.
\]

\begin{enumerate}
\def\labelenumi{(\alph{enumi})}
\setcounter{enumi}{3}
\tightlist
\item
  Let I be a well-ordered set and \((\lambda_{t})_{i\in I}\) a family of order-types indexed by I. The order-type of the lexicographic product of the family of sets ordered by the \(\lambda_{t} (t\in I)\) is called the ordinal product of the family \((\lambda_{t})_{i\in I}\) and is denoted by \(\mathsf{P}_{i\in I}\lambda_{t}\) . If \((E_{t})_{i\in I}\) is a family of ordered sets, the order-type of the lexicographic product of the family \((E_{t})_{i\in I}\) is \(\mathsf{P}_{i\in I}\operatorname{Ord}(E_{t})\) . If I is the ordinal sum of a family of well-ordered sets \((J_{x})_{x\in K}\) , indexed by a well-ordered set \(K\) , show that
\end{enumerate}

\[
\underset {x \in \mathbf {K}} {\mathsf {P}} \left(\underset {i \in J _ {x}} {\mathsf {P}} \lambda_ {i}\right) = \underset {i \in I} {\mathsf {P}} \lambda_ {i}.
\]

\begin{enumerate}
\def\labelenumi{(\alph{enumi})}
\setcounter{enumi}{4}
\item
  We denote by \(\lambda + \mu\) (resp. \(\mu\lambda\) ) the ordinal sum (resp. ordinal product) of the family \((\xi_{i})_{i \in J}\) , where \(J = \{\alpha, \beta\}\) is a set with two distinct elements, ordered by the relation whose graph is \(\{(\alpha, \alpha), (\alpha, \beta), (\beta, \beta)\}\) , and where \(\xi_{\alpha} = \lambda\) , \(\xi_{\beta} = \mu\) . Show that if I is a well-ordered set of order-type \(\lambda\) , and if \((\mu_{i})_{i \in I}\) is a family of order-types such that \(\mu_{i} = \mu\) for each \(i \in I\) , then \(\sum_{i \in I} \mu_{i} = \mu\lambda\) . We have \((\lambda + \mu) + \nu = \lambda + (\mu + \nu)\) , \((\lambda\mu)\nu = \lambda(\mu\nu)\) , and \(\lambda(\mu + \nu) = \lambda\mu + \lambda\nu\) (but in general \(\lambda + \mu \neq \mu + \lambda\) , \(\lambda\mu \neq \mu\lambda\) , \((\lambda + \mu)\nu \neq \lambda\nu + \mu\nu\) ).
\item
  Let \((\lambda_{t})_{t\in I}\) and \((\mu_{t})_{t\in I}\) be two families of order-types indexed by the same ordered set \(I\) . Show that, if \(\lambda_{t} \prec \mu_{t}\) for each \(t \in I\) , then \(\sum_{i \in I} \lambda_{t} \prec \sum_{i \in I} \mu_{t}\) and (if \(I\) is well-ordered) \(\mathsf{P}_{t} \lambda_{t} \prec \mathsf{P}_{t \in I} \mu_{t}\) . If \(J\) is a subset of \(I\) , show that \(\sum_{i \in J} \lambda_{t} \prec \sum_{i \in I} \lambda_{t}\) and (if \(I\) is well-ordered and the \(\lambda_{t}\) are non-empty) \(\mathsf{P}_{t \in J} \lambda_{t} \prec \mathsf{P}_{t \in I} \lambda_{t}\) .
\item
  Let \(\lambda^{*}\) denote the order-type of the set ordered by the opposite of the ordering \(\lambda\) . Then we have
\end{enumerate}

\[
(\lambda^ {*}) ^ {*} = \lambda \quad \text { and } \quad \left(\sum_ {i \in I} \lambda_ {i}\right) ^ {*} = \sum_ {i \in I ^ {*}} \lambda_ {i},
\]

where I* denotes the set I endowed with the opposite of the ordering given on I.

\begin{enumerate}
\def\labelenumi{\arabic{enumi}.}
\setcounter{enumi}{13}
\tightlist
\item
  An ordinal is the order-type of a well-ordered set (Exercise 13).
\end{enumerate}

\begin{enumerate}
\def\labelenumi{(\alph{enumi})}
\tightlist
\item
  Show that, if \((\lambda_{i})_{i\in I}\) is a family of ordinals indexed by a well-ordered set \(I\) , then the ordinal sum \(\sum_{i\in I}\lambda_{i}\) is an ordinal; * and that if, moreover, \(I\) is finite, the ordinal product \(\mathsf{P}_{i\in I}\lambda_{i}\) is an ordinal (Exercise 11).* The order-type of the empty set is denoted by 0, and that of a set with one element by 1 (by abuse of language, cf.~§3). Show that
\end{enumerate}

\[
\alpha + 0 = 0 + \alpha = \alpha \quad \text { and } \quad \alpha . 1 = 1. \alpha = \alpha
\]

for every ordinal \(\alpha\) .

\begin{enumerate}
\def\labelenumi{(\alph{enumi})}
\setcounter{enumi}{1}
\item
  Show that the relation `` \(\lambda\) is an ordinal and \(\mu\) is an ordinal and \(\lambda \prec \mu\) '' is a well-ordering relation, denoted by \(\lambda \leqslant \mu\) . (Note that, if \(\lambda\) and \(\mu\) are ordinals, the relation \(\lambda \prec \mu\) is equivalent to `` \(\lambda\) is equal to the order-type of a segment of \(\mu\) '' (no. 5, Theorem 3, Corollary 3): given a family \((\lambda_{i})_{i \in I}\) of ordinals, consider a well-ordering on \(I\) and take the ordinal sum of the family of sets ordered by the \(\lambda_{i}\) ; finally, use Proposition 2 of no. 1.)
\item
  Let \(\alpha\) be an ordinal. Show that the relation `` \(\xi\) is an ordinal and \(\xi \leqslant \alpha\) '' is collectivizing in \(\xi\) , and that the set \(O_{\alpha}\) of ordinals \(< \alpha\) is a well-ordered set such that \(\operatorname{Ord}(O_{\alpha}) = \alpha\) . We shall often identify \(O_{\alpha}\) with \(\alpha\) .
\item
  Show that for every family of ordinals \((\xi_{i})_{i\in I}\) there exists a unique ordinal \(\alpha\) such that the relation `` \(\lambda\) is an ordinal and \(\xi_{i}\leqslant\lambda\) for all \(i\in I\) '' is equivalent to \(\alpha\leqslant\lambda\) . By abuse of language, \(\alpha\) is called the least upper bound of the family of ordinals \((\xi_{i})_{i\in I}\) , and we write \(\alpha=\sup_{i\in I}\xi_{i}\)
\end{enumerate}

(it is the greatest element of the union of \(\{\alpha\}\) and the set of the \(\xi_{i}\) ). The least upper bound of the set of ordinals \(\xi < \alpha\) is either \(\alpha\) or an ordinal \(\beta\) such that \(\alpha = \beta + 1\) ; in the latter case, \(\beta\) is said to be the predecessor of \(\alpha\) .

\begin{enumerate}
\def\labelenumi{\arabic{enumi}.}
\setcounter{enumi}{14}
\tightlist
\item
  \begin{enumerate}
  \def\labelenumii{(\alph{enumii})}
  \tightlist
  \item
    Let \(\alpha\) and \(\beta\) be two ordinals. Show that the inequality \(\alpha < \beta\) is equivalent to \(\alpha + 1 \leqslant \beta\) , and that it implies the inequalities \(\xi + \alpha < \xi + \beta\) , \(\alpha + \xi \leqslant \beta + \xi\) , \(\alpha\xi \leqslant \beta\xi\) for all ordinals \(\xi\) , and \(\xi\alpha < \xi\beta\) if \(\xi > 0\) .
  \end{enumerate}
\end{enumerate}

\begin{enumerate}
\def\labelenumi{(\alph{enumi})}
\setcounter{enumi}{1}
\item
  Deduce from (a) that there exists no set to which every ordinal belongs (use Exercise 14 (d)).
\item
  Let \(\alpha, \beta, \mu\) be three ordinals. Show that each of the relations \(\mu + \alpha < \mu + \beta\) , \(\alpha + \mu < \beta + \mu\) implies \(\alpha < \beta\) ; and that each of the relations \(\mu \alpha < \mu \beta\) , \(\alpha \mu < \beta \mu\) implies \(\alpha < \beta\) provided that \(\mu > 0\) .
\item
  Show that the relation \(\mu + \alpha = \mu + \beta\) implies \(\alpha = \beta\) , and that \(\mu \alpha = \mu \beta\) implies \(\alpha = \beta\) provided that \(\mu > 0\) .
\item
  Two ordinals \(\alpha\) and \(\beta\) are such that \(\alpha \leqslant \beta\) if and only if there exists an ordinal \(\xi\) such that \(\beta = \alpha + \xi\) . This ordinal \(\xi\) is then unique and is such that \(\xi \leqslant \beta\) ; it is written \((- \alpha) + \beta\) .
\item
  Let \(\alpha, \beta, \zeta\) be three ordinals such that \(\zeta < \alpha\beta\) . Show that there exist two ordinals \(\xi, \eta\) such that \(\zeta = \alpha\eta + \xi\) and \(\xi < \alpha, \eta < \beta\) (cf.~no. 5, Theorem 3, Corollary 3). Moreover, \(\xi\) and \(\eta\) are uniquely determined by these conditions.
\end{enumerate}

\begin{enumerate}
\def\labelenumi{\arabic{enumi}.}
\setcounter{enumi}{15}
\tightlist
\item
  An ordinal \(\rho > 0\) is said to be indecomposable if there exists no pair of ordinals \(\xi, \eta\) such that \(\xi < \rho, \eta < \rho\) , and \(\xi + \eta = \rho\) .
\end{enumerate}

\begin{enumerate}
\def\labelenumi{(\alph{enumi})}
\item
  An ordinal \(\rho\) is indecomposable if and only if \(\xi + \rho = \rho\) for every ordinal \(\xi\) such that \(\xi < \rho\) .
\item
  If \(\rho > 1\) is an indecomposable ordinal and if \(\alpha\) is any ordinal \(>0\) , then \(\alpha\rho\) is indecomposable, and conversely (use Exercise 15 (f)).
\item
  If \(\rho\) is indecomposable and if \(0 < \alpha < \rho\) , then \(\rho = \alpha\xi\) , where \(\xi\) is an indecomposable ordinal (use Exercise 15 (f)).
\item
  Let \(\alpha\) be an ordinal \(>0\) . Show that there exists a greatest indecomposable ordinal among the indecomposable ordinals \(\leqslant \alpha\) (consider the decompositions \(\alpha = \rho + \xi\) , where \(\rho\) is indecomposable).
\item
  If \(\mathbf{E}\) is any set of indecomposable ordinals, deduce from (d) that the least upper bound of \(\mathbf{E}\) (Exercise 14 (d)) is an indecomposable ordinal.
\end{enumerate}

\begin{enumerate}
\def\labelenumi{\arabic{enumi}.}
\setcounter{enumi}{16}
\tightlist
\item
  Given an ordinal \(\alpha_0\) , a term \(f(\xi)\) is said to be an ordinal functional symbol (with respect to \(\xi\) ) defined for \(\xi \geqslant \alpha_0\) if the relation `` \(\xi\) is an ordinal and \(\xi \geqslant \alpha_0\) '' implies the relation `` \(f(\xi)\) is an ordinal''; \(f(\xi)\) is said to be normal if the relation \(\alpha_0 \leqslant \xi < \eta\) implies \(f(\xi) < f(\eta)\) and if, for each family \((\xi_t)_{t \in I}\) of ordinals \(\geqslant \alpha_0\) , we have \(\sup_{t \in I} f(\xi_t) = f(\sup_{t \in I} \xi_t)\) (cf.~Exercise 14 (d)).
\end{enumerate}

\begin{enumerate}
\def\labelenumi{(\alph{enumi})}
\item
  Show that, for each ordinal \(\alpha > 0\) , \(\alpha + \xi\) and \(\alpha\xi\) are normal functional symbols defined for \(\xi \geqslant 0\) (use Exercise 15 (f)).
\item
  Let \(w(\xi)\) be an ordinal functional symbol defined for \(\xi \geqslant \alpha_0\) such that \(w(\xi) \geqslant \xi\) and such that \(\alpha_0 \leqslant \xi < \eta\) implies \(w(\xi) < w(\eta)\) . Also let \(g(\xi, \eta)\) be a term such that the relation `` \(\xi\) and \(\eta\) are ordinals \(\geqslant \alpha_0\) '' implies the relation `` \(g(\xi, \eta)\) is an ordinal such that \(g(\xi, \eta) > \xi\) ''.
\end{enumerate}

Define a term \(f(\xi, \eta)\) with the following properties: (1) for each ordinal \(\xi \geqslant \alpha_0\) , \(f(\xi, 1) = w(\xi)\) ; (2) for each ordinal \(\xi \geqslant \alpha_0\) and each ordinal \(\eta > 1\) , \(f(\xi, \eta) = \sup_{0 < \zeta < \eta} g(f(\xi, \zeta), \xi)\) (use criterion C60 of no. 2). Show that if \(f_1(\xi, \eta)\) is another term with these two properties, then \(f(\xi, \eta) = f_1(\xi, \eta)\) for all \(\xi \geqslant \alpha_0\) and all \(\eta \geqslant 1\) . Prove that, for each ordinal \(\xi \geqslant \alpha_0\) , \(f(\xi, \eta)\) is a normal functional symbol with respect to \(\eta\) (defined for all \(\eta \geqslant 1\) ). Show that \(f(\xi, \eta) \geqslant \xi\) for all \(\eta \geqslant 1\) and \(\xi \geqslant \alpha_0\) , and that \(f(\xi, \eta) \geqslant \eta\) for all \(\xi \geqslant \sup (\alpha_0, 1)\) and \(\eta \geqslant 1\) . Furthermore, for each pair \((\alpha, \beta)\) of ordinals such that \(\alpha > 0\) , \(\alpha \geqslant \alpha_0\) , and \(\beta \geqslant w(\alpha)\) there exists a unique ordinal \(\xi\) such that

\[
f (\alpha , \xi) \leqslant \beta <   f (\alpha , \xi + 1),
\]

and we have \(\xi \leqslant \beta\) .

\begin{enumerate}
\def\labelenumi{(\alph{enumi})}
\setcounter{enumi}{2}
\item
  If we take \(\alpha_0 = 0\) , \(w(\xi) = \xi + 1\) , \(g(\xi, \eta) = \xi + 1\) , then \(f(\xi, \eta) = \xi + \eta\) . If we take \(\alpha_0 = 1\) , \(w(\xi) = \xi\) , \(g(\xi, \eta) = \xi + \eta\) , then \(f(\xi, \eta) = \xi \eta\) .
\item
  Show that if the relations \(\alpha_0 \leqslant \xi \leqslant \xi'\) , \(\alpha_0 \leqslant \eta \leqslant \eta'\) imply \(g(\xi, \eta) \leqslant g(\xi', \eta')\) , then the relations \(\alpha_0 \leqslant \xi \leqslant \xi'\) , \(1 \leqslant \eta \leqslant \eta'\) imply \(f(\xi, \eta) \leqslant f(\xi', \eta')\) . If the relations \(\alpha_0 \leqslant \xi \leqslant \xi'\) , \(\alpha_0 \leqslant \eta < \eta'\) imply \(g(\xi, \eta) < g(\xi, \eta')\) and \(g(\xi, \eta) \leqslant g(\xi', \eta)\) , then the relations \(\alpha_0 \leqslant \xi < \xi'\) and \(\eta \geqslant 0\) imply \(f(\xi, \eta + 1) < f(\xi', \eta + 1)\) .
\item
  Suppose that \(w(\xi) = \xi\) and that the relations \(\alpha_0 \leqslant \xi \leqslant \xi'\) , \(\alpha_0 \leqslant \eta < \eta'\) imply \(g(\xi, \eta) < g(\xi, \eta')\) and \(g(\xi, \eta) \leqslant g(\xi', \eta)\) . Suppose, moreover, that for each \(\xi \geqslant \alpha_0\) , \(g(\xi, \eta)\) is a normal functional symbol with respect to \(\eta\) (defined for \(\eta \geqslant \alpha_0\) ), and that, whenever \(\xi \geqslant \alpha_0\) , \(\eta \geqslant \alpha_0\) , and \(\zeta \geqslant \alpha_0\) , we have the associativity relation
\end{enumerate}

\[
g (g (\xi , \eta), \zeta) = g (\xi , g (\eta , \zeta)).
\]

Show that, if \(\xi \geqslant \alpha_0\) , \(\eta \geqslant 1\) , and \(\zeta \geqslant 1\) , then

\[
g (f (\xi , \eta), f (\xi , \zeta)) = f (\xi , \eta + \zeta)
\]

(``distributivity'' of \(g\) with respect to \(f\) ) and

\[
f (f (\xi , \eta), \zeta) = f (\xi , \eta \zeta)
\]

(``associativity'' of \(f\) ).

¶ 18. In the definition procedure defined in Exercise 17 (b), take \(\alpha_{0}=1+1\) (denoted by 2 by abuse of language),

\[
w (\xi) = \xi , \quad g (\xi , \eta) = \xi \eta .
\]

Denote \(f(\xi, \eta)\) by \(\xi^{\eta}\) and define \(\alpha^{0}\) to be 1 for all ordinals \(\alpha\) . Also define \(0^{\beta}\) to be 0 and \(1^{\beta}\) to be 1 for all ordinals \(\beta \geqslant 1\) .

\begin{enumerate}
\def\labelenumi{(\alph{enumi})}
\item
  Show that if \(\alpha > 1\) and \(\beta < \beta'\) , we have \(\alpha^{\beta} < \alpha^{\beta'}\) and that, for each ordinal \(\alpha > 1\) , \(\alpha^{\xi}\) is a normal functional symbol with respect to \(\xi\) . Moreover, if \(0 < \alpha \leqslant \alpha'\) , we have \(\alpha^{\beta} \leqslant \alpha'^{\beta}\) .
\item
  Show that \(\alpha^{\xi}.\alpha^{\eta} = \alpha^{\xi +\eta}\) and \((\alpha^{\xi})^{\eta} = \alpha^{\xi \eta}\) .
\item
  Show that, if \(\alpha \geqslant 2\) and \(\beta \geqslant 1\) , \(\alpha^{\beta} \geqslant \alpha \beta\) .
\item
  For each pair of ordinals \(\beta \geqslant 1\) and \(\alpha \geqslant 2\) , there exist three ordinals \(\xi, \gamma, \delta\) such that \(\beta = \alpha^{\xi}\gamma + \delta\) , where \(0 < \gamma < \alpha\) and \(\delta < \alpha^{\xi}\) , and these three ordinals are uniquely determined by these conditions.
\end{enumerate}

* 19. Let \(\alpha\) and \(\beta\) be two ordinals, and let E, F be two well-ordered sets such that \(\operatorname{Ord}(E) = \alpha\) and \(\operatorname{Ord}(F) = \beta\) . In the set \(E^F\) of mappings of F into E, consider the subset G of mappings \(g\) such that \(g(y)\) is equal to the least element of E for all but a finite number of elements \(y \in F\) . If \(F^*\) is the ordered set obtained by endowing F with the opposite order, show that G is a connected component with respect to the relation `` \(x\) and \(y\) are comparable'' (Chapter II, §6, Exercise 10) in the product \(E^{F*}\) endowed with the ordering defined in Exercise 12, and show that G is well-ordered. Furthermore, prove that \(\operatorname{Ord}(G) = \alpha^\beta\) (use the uniqueness property of Exercise 17 (b)).

\begin{enumerate}
\def\labelenumi{\arabic{enumi}.}
\setcounter{enumi}{19}
\tightlist
\item
  A set X is said to be transitive if the relation \(x \in \mathbf{X}\) implies \(x \subset \mathbf{X}\) .
\end{enumerate}

\begin{enumerate}
\def\labelenumi{(\alph{enumi})}
\item
  If Y is a transitive set, then so is \(Y \cup \{Y\}\) . If \((Y_{t})_{t \in I}\) is a family of transitive sets, then \(\bigcup_{i \in I} Y_{t}\) and \(\bigcap_{i \in I} Y_{t}\) are transitive.
\item
  A set X is a pseudo-ordinal if every transitive set Y such that \(Y \subset X\) and \(Y \neq X\) is an element of X. A set S is said to be decent if the relation \(x \in S\) implies \(x \notin x\) . Show that every pseudo-ordinal is transitive and decent (consider the union of the decent transitive subsets of X, and use (a)). If X is a pseudo-ordinal, then so is \(X \cup \{X\}\) .
\item
  Let X be a transitive set, and suppose that each \(x \in X\) is a pseudo-ordinal. Then X is a pseudo-ordinal (note that, for each \(x \in X\) , \(x \cup \{x\}\) is a pseudo-ordinal contained in X).
\item
  Show that \(\phi\) is a pseudo-ordinal and that every element of a pseudo-ordinal X is a pseudo-ordinal. (Consider the union of the transitive subsets of X whose elements are pseudo-ordinals.)
\item
  If \((\mathbf{X}_t)_{t\in \mathbf{I}}\) is a family of pseudo-ordinals, then \(\bigcap_{t\in \mathbf{I}}\mathbf{X}_t\) is the least element of this family (with respect to the relation of inclusion). (Use (b).) Deduce that, if E is a pseudo-ordinal, the relation \(x\subset y\) between elements \(x,y\) of E is a well-ordering relation.
\item
  Show that for each ordinal \(\alpha\) there exists a unique pseudo-ordinal \(\mathbf{E}_{\alpha}\) such that \(\operatorname{Ord}(\mathbf{E}_{\alpha}) = \alpha\) (use (e) and criterion C60). In particular, the pseudo-ordinals whose order-types are \(0, 1, 2 = 1 + 1\) , and \(3 = 2 + 1\) are respectively
\end{enumerate}

\[
\phi , \quad \{\phi \}, \quad \{\phi , \{\phi \} \}, \quad \{\phi , \{\phi \}, \{\phi , \{\phi \} \} \}.
\]

